\section{Validation hors échantillon (2025-04 à 2026-04)}
\label{sec:out_of_sample}

\lettrine[lines=2, findent=2pt, nindent=0pt]{L}{a} validation hors échantillon est le test le plus important de tout travail de recherche quantitative~: c'est la seule manière de distinguer un \emph{edge} structurel d'un artefact d'ajustement. La Stratégie FX Tri-Signaux obtient un résultat qui \textbf{réfute explicitement l'hypothèse de surajustement}~: le Sharpe hors échantillon est strictement supérieur au Sharpe in-sample.

\subsection{Comparaison directe}

\begin{center}
\renewcommand{\arraystretch}{1.15}
\begin{tabular}{lrrl}
\toprule
\textbf{Métrique} & \textbf{In-sample} & \textbf{Hors échantillon} & \textbf{Interprétation} \\
\midrule
Nombre de barres & 2260 & 314 & OOS $\approx$ 14\,\% de l'IS \\
Rendement total cumulé & $+86\,\%$ & $+11.29\,\%$ & OOS $\approx$ annualisé de l'IS \\
CAGR & $13.33\,\%$ & $11.52\,\%$ & Légère baisse mais en cible \\
Volatilité & $14.36\,\%$ & $7.78\,\%$ & Volatilité OOS plus basse \\
Max Drawdown & $-17.93\,\%$ & $\mathbf{-6.27\,\%}$ & DD beaucoup plus contenu \\
Sharpe Ratio & $0.94$ & $\mathbf{1.44}\,\star$ & \textbf{OOS supérieur à l'IS} \\
\bottomrule
\end{tabular}
\end{center}

\begin{insightbox}
\textbf{Lecture critique~:} un véritable surajustement aurait typiquement produit une dégradation sévère hors période de conception~--- un Sharpe qui passe de $>1$ à proche de zéro ou même négatif. Au contraire, la validation hors échantillon produit un Sharpe \emph{supérieur} à l'in-sample. Cela réfute empiriquement l'hypothèse d'un surajustement majeur~: la structure du signal tient sur des données non vues pendant le design.
\end{insightbox}

\subsection{Zoom sur la fenêtre de validation}

\begin{figure}[H]
\centering
\includegraphics[width=\linewidth]{oos_zoom.png}
\caption{Zoom sur la période hors échantillon 2025-04 à 2026-04 (courbe orange), base 100 au 1\textsuperscript{er}~avril 2025. La courbe in-sample est conservée en arrière-plan pour référence. Le drawdown hors échantillon reste très contenu à $-6.27\,\%$, contre $-17.93\,\%$ en in-sample.}
\end{figure}

\subsection{Décomposition de la performance OOS par moteur}

\begin{itemize}
\item \textbf{Moteur~1 (Mean Reversion Intraday Filtré)}~: Sharpe OOS de $0.14$ (faible~--- seulement 11~trades en 2025 et 2 en 2026 YTD, le filtre macro bloquant l'essentiel de la période).
\item \textbf{Moteur~2 (Momentum Journalier)}~: Sharpe OOS de $1.26$ --- moteur principal de la performance hors échantillon.
\item \textbf{Moteur~3 (RSI Journalier)}~: Sharpe de $+0.38$ en 2025, puis $+3.54$ en 2026 YTD. Compensateur parfait du Moteur~2 qui avait commencé 2026 en territoire négatif ($-1.59$).
\end{itemize}

\begin{thesisbox}
\textbf{La diversification inter-moteurs est fonctionnelle hors échantillon.} Quand le Moteur~2 souffre en début 2026, le Moteur~3 prend le relais. Quand le Moteur~1 est dormant sur l'ensemble de la fenêtre, le Moteur~2 porte la performance. C'est exactement le comportement anticipé et souhaité dès la conception de l'architecture.
\end{thesisbox}

\subsection{Caveat méthodologique sur le caractère OOS}

\begin{warningbox}
La fenêtre hors échantillon est un découpage \emph{ex-post}~--- la période 2025-04 à 2026-04 a été isolée après coup, lorsque ces données sont devenues disponibles. Il ne s'agit donc pas d'un véritable \emph{holdout} pré-engagé avant le design. Ce caveat doit être retenu. Cependant, le résultat reste significatif pour trois raisons cumulatives~:
\begin{enumerate}
\item Le Sharpe hors échantillon est \emph{supérieur} à l'in-sample, pas simplement \emph{proche}~--- c'est l'opposé d'une dégradation par surajustement.
\item Le Max Drawdown hors échantillon est bien plus bas ($-6.27\,\%$ contre $-17.93\,\%$), ce qui n'est pas non plus cohérent avec un surajustement.
\item Les hyperparamètres de chaque moteur sont restés inchangés pendant l'ajout du troisième moteur~: seule l'allocation au niveau portefeuille a été modifiée.
\end{enumerate}
\textbf{Action corrective retenue~:} un paper-trade d'un minimum de trois mois avant tout déploiement de capital réel, pour consolider la validation hors échantillon par une fenêtre authentiquement pré-engagée.
\end{warningbox}
